% page 484 du fac-simile (image p-483 du scan)
% bloc 157.5 x 233.3 mm ; marge gauche 8.0 mm
\pgc{-12.700mm}{0.000mm}{
\l{\cel{3}476}
\l{}
\l{\sou{quartana}. (\sou{patol.}) Dicesas pri febro intermitanta qua ri-}
\l{\cel{3}aparas singla-quardie. - DEFIS.}
\l{}
\l{\sou{quartero}. Segmento administrala di urbo. - DEFIRS.}
\l{}
\l{\sou{quartermastro}. (\sou{navig.}) Grado unesma super olta di navano,}
\l{\cel{3}che la navistari di la militistaro marala. (Lu korespon-}
\l{\cel{3}das a la grado \textquotedbl{}kaporalo\textquotedbl{} di la militistaro terala).-}
\l{\cel{3}DEFIRS.}
\l{}
\l{\sou{quarterono}. Homo qua devenis de la kopulaco da pelblanko kun}
\l{\cel{3}mestico. - DEFIRS.}
\l{}
\l{\sou{quartebo}.\cel{2}(muziko). I. Muzikajo kompozita por quar voci o}
\l{\cel{3}por quar instrumenti (du violini, alto e violoncelo). -}
\l{\cel{3}II. (en\cel{1}\sou{orkestro}).Ensemblo ek la violini unesma e duesma,}
\l{\cel{3}alti, violonceli (o kontrabasi) qui formacas quar parti.}
\l{\cel{3}DEFIRS.}
\l{}
\l{\sou{quartiko}. (\sou{geom.)}\cel{2}Kurvo di la klaso quaresma. - EFIS.}
\l{}
\l{\sou{quasio.(bot.}) Kortico de la arbusto L.\cel{1}\sou{quassia amara}, de-}
\l{\cel{3}tranchita aden mikra spani, qua saporas tre bitre, ed}
\l{\cel{3}uzata en medicino kom tonigivo pos macereso (5 grami}
\l{\cel{3}singla-litre), sive en tintivo alkoholoza. DEFIS.}
\l{}
\l{\sou{quaternara}. I. (\sou{mat.)}\cel{2}En qua quar unaji simpla formacas}
\l{\cel{3}unajo ek la klaso quaresma, e tale pluse. - II. (\sou{geol.})}
\l{\cel{3}\sou{Strato quaternara}\cel{1}: la maxim recenta ek la strati di for-}
\l{\cel{3}maceso sedimentala.- DEFIS.}
\l{}
\l{\sou{quaterniono}. (\sou{matem.)}\cel{2}Nomo donita da Hamilton ad ula expre-}
\l{\cel{3}suri imaginara por qui lu kreis kalkulo-speco partiku-}
\l{\cel{3}lara qua posibligas la solvo facila di la problemi qui}
\l{\cel{3}koncernas la geometrio tri-dimensiona. - DEFIS.}
\l{}
\l{\sou{quatreno}. Stanco, parto di soneto, o mikra poeziajo quar-}
\l{\cel{3}versa. - DEFIS.}
\l{}
\l{\sou{quaza}. Simila til ula grado, ma ne tote tala. (Indikas nur}
\l{\cel{3}semblo o simileso, e ke la du objekti interdiferas per}
\l{\cel{3}nur poko). - DEFIS.}
\l{}
\l{\sou{quecho}. (\sou{bot.}) Sorto di grosa pruno oblonga, qua su repro-}
\l{\cel{3}duktas per semuro. - L.\cel{1}\sou{prunus domestica}. - - F.}
\l{}
\l{\sou{querar}.\cel{2}(\sou{trans.})\cel{2}Serchar, komisite od intencante adportor}
\l{\cel{3}od venigor. - FI.}
\l{}
\l{\sou{quercitrono}. (\sou{bot.)}\cel{1}Granda querko sempre-verda de la nordo-}
\l{\cel{3}regioni di Amerika, di qua la kortico donas materio tintala}
\l{\cel{3}flava. L.\cel{1}\sou{quercus tinctoria.}\cel{2}DEFI.}
}
